\section{Thiết kế hai luồng vận hành chính}

Phân hệ RAG được chia thành hai luồng vận hành chính: \textbf{Indexing Flow} và \textbf{Query Flow}. Indexing Flow hoạt động ngoại tuyến hoặc định kỳ để đưa tài liệu mới vào kho tri thức. Query Flow hoạt động trực tuyến khi người dùng đặt câu hỏi.

\subsection{Indexing Flow - luồng sản xuất tri thức}

Indexing Flow biến tài liệu tuyển sinh thô thành tri thức có thể truy hồi. Đầu vào là PDF, DOCX, HTML, website hoặc văn bản nhập tay; đầu ra là các chunk đã gắn metadata và embedding.

Các bước chính:
\begin{itemize}
    \item Tải tài liệu lên, kiểm tra định dạng và dung lượng.
    \item Tính checksum để phát hiện trùng lặp.
    \item Trích xuất văn bản (kèm OCR nếu là ảnh/PDF scan).
    \item Làm sạch, chuẩn hóa và chia chunk.
    \item Gán metadata, sinh embedding, lưu vào các lớp lưu trữ.
    \item Cập nhật trạng thái tài liệu thành đã lập chỉ mục.
\end{itemize}

\begin{figure}[H]
    \centering
    \includegraphics[width=0.95\textwidth]{\detokenize{content/source/hình 5.2 Sơ đồ Indexing Flow..png}}
    \caption{Sơ đồ Indexing Flow}
    \label{fig:indexing-flow}
\end{figure}

\subsection{Query Flow - luồng tiêu thụ tri thức}

Query Flow xử lý câu hỏi tự nhiên của người dùng và trả về câu trả lời có trích dẫn hoặc thông báo không đủ dữ liệu.

Các bước chính:
\begin{itemize}
    \item Người dùng gửi câu hỏi qua chatbot.
    \item Rasa phân tích intent/entity, sau đó gọi Custom Action.
    \item Custom Action gửi truy vấn đến RAG Query API.
    \item RAG chuẩn hóa câu hỏi, truy hồi hybrid (Vector Search + BM25), reranking.
    \item Context Builder chọn bằng chứng, Prompt Assembly tạo prompt.
    \item LLM sinh câu trả lời, Citation Module gắn nguồn.
    \item Trả kết quả về Rasa để hiển thị cho người dùng.
\end{itemize}

\begin{figure}[H]
    \centering
    \includegraphics[width=0.95\textwidth]{\detokenize{content/source/hình 5.3 Sơ đồ Query Flow..png}}
    \caption{Sơ đồ Query Flow}
    \label{fig:query-flow}
\end{figure}
